\section{Liniowe operatory w przestrzeni unitarnej}

Rozpatrzmy $n$-wymiarową przestrzeń unitarną $V$. Niech będzie dany dowolny operator liniowy $A$ w tej przestrzeni.

\begin{definition}
    Operator liniowy $A^*$ nazywa się \textbf{operatorem sprzężonym do operatora} $A$, jeśli dla dowolnych dwóch wektorów $x, y \in V$:
    \begin{equation}\label{eq:def_conjugated_operator}
        \inner{Ax}{y} = \inner{x}{A^*y}
    \end{equation}
\end{definition}
\begin{lemma}\label{lemma:existing_of_ortonormal_base}
W każdej skończonie wymiarowej przestrzeni unitarnej $V$ można wybrać bazę ortonormalną.
\end{lemma}   
\begin{proof}
    Dowolną bazę $V$ można zortogonalizować metodą Grama–Schmidta \cite[s. 140--142]{mostowski1977}, a następnie znormalizować otrzymane wektory, co daje bazę ortonormalną.
\end{proof}

\begin{theorem}\label{theorem:existing_of_conjugated_operator}
Dla każdego liniowego operatora $A$ w przestrzeni unitarnej $V$ zawsze istnieje operator sprzężony $A^*$ i to tylko jeden.
\end{theorem}
\begin{proof}
\textbf{Istnienie.} Niech $\mathbf{e}_1, \ldots, \mathbf{e}_n$ jest bazą ortonormalną przestrzeni $V$. Dowolny wektor $y \in V$ w bazie ortonormalnej da się zapisać w postaci:
$$ y = \sum_{k=1}^{n}{\inner{y}{\mathbf{e}_k}\mathbf{e}_k} $$
    Wtedy dla wektora $y$ i szukanego operatora $A^*$:
\begin{equation}
        A^*y = \sum_{k=1}^{n}{\inner{A^*y}{\mathbf{e}_k}\mathbf{e}_k}
\end{equation}
\begin{equation}
        A^*y = \sum_{k=1}^{n}{\inner{y}{A\mathbf{e}_k}\mathbf{e}_k}
\end{equation}

Operator ten jest liniowy, ponieważ iloczyn skalarny jest liniowy.
Pokażemy, że tak zdefiniowany operator spełnia warunek \eqref{eq:def_conjugated_operator}. Niech $x, y \in V$ będą dowolne. Zapiszmy $x$ w bazie ortonormalnej: $x = \sum_{k=1}^n \inner{x}{\mathbf{e}_k}\mathbf{e}_k$. Wówczas
\begin{equation}
    \inner{x}{A^*y} = \Big\langle \sum_{k=1}^n \inner{x}{\mathbf{e}_k}\mathbf{e}_k,\ \sum_{i=1}^n \inner{y}{A\mathbf{e}_i}\mathbf{e}_i \Big\rangle = \sum_{k=1}^n \inner{x}{\mathbf{e}_k}\overline{\inner{y}{A\mathbf{e}_k}},
\end{equation}
gdzie skorzystaliśmy z ortonormalności $\base{e}$. 
Z drugiej strony
\begin{equation}
    \inner{Ax}{y} = \inner{
        A\Big(\sum_{k=1}^n \inner{x}{\mathbf{e}_k}\mathbf{e}_k\Big)
    }{y} = \sum_{k=1}^n \inner{x}{\mathbf{e}_k}\inner{A\mathbf{e}_k}{y} = \sum_{k=1}^n \inner{x}{\mathbf{e}_k}\overline{\inner{y}{A\mathbf{e}_k}},
\end{equation}
gdzie ostatnia równość wynika z symetrii hermitowskiej iloczynu skalarnego w przestrzeni unitarnej.
Prawe strony obu równości są identyczne, więc $\inner{Ax}{y} = \inner{x}{A^*y}$ dla wszystkich $x,y \in V$, co dowodzi istnienia $A^*$.

\textbf{Jednoznaczność.} Załóżmy, że operatory $B_1, B_2$ oba spełniają warunek \eqref{eq:def_conjugated_operator}, tzn. $\inner{Ax}{y} = \inner{x}{B_1y} = \inner{x}{B_2y}$ dla wszystkich $x,y \in V$. Wtedy dla dowolnych $x, y \in V$:
\begin{equation}
    \inner{x}{B_1y - B_2y} = 0.
\end{equation}
Ponieważ $x \in V$ było dowolne, a iloczyn skalarny w $V$ jest niezdegenerowany (tzn. $\inner{x}{z}=0$ dla wszystkich $x$ różnych od zera implikuje $z=0$), otrzymujemy $B_1y - B_2y = 0$ dla każdego $y \in V$, czyli $B_1 = B_2$.
\end{proof}

\begin{theorem}
    W przestrzeni unitarnej $V$ z bazą ortonormalną, sprzężonemu operatorowi $A^*$ odpowiada macierz sprzężona (tzn. $A^* = \overline{A^T}$).
\end{theorem}
\begin{proof}
    Niech $A$ będzie operatorem liniowym w przestrzeni unitarnej $V$ z  ortonormalną bazą $\base{e}$ (Lemat \ref{lemma:existing_of_ortonormal_base}) któremu odpowiada macierz $(a_{ik})$, gdzie $i, k = 1, \ldots, n$.

    Zastosujmy do dowolnego wektora $\mathbf{e}_k$ z bazy $\base{e}$ operator liniowy $A$:
    \begin{equation}\label{eq:local_Ae_def}
        A\mathbf{e}_k = \sum_{i=1}^{n}{a_{ik}\mathbf{e}_i} = a_{1k}\mathbf{e}_1 + \ldots + a_{ik}\mathbf{e}_i + \ldots + a_{nk}\mathbf{e}_n
    \end{equation}
    Pomnożymy skalarnie obie strony \eqref{eq:local_Ae_def} przez dowolny wektor bazowy $\mathbf{e}_i$, gdzie $i \in \{1, \ldots, n\}$:
    $$\inner{A\mathbf{e}_k}{\mathbf{e}_i} = \inner{
        a_{1k}\mathbf{e}_1 + \ldots + a_{ik}\mathbf{e}_i + \ldots + a_{nk}\mathbf{e}_n
    }{\mathbf{e}_i}$$
    $$\inner{A\mathbf{e}_k}{\mathbf{e}_i} = a_{1k}\inner{\mathbf{e}_1}{\mathbf{e}_i} + \ldots + a_{ik}\inner{\mathbf{e}_i}{\mathbf{e}_i} + \ldots + a_{nk}\inner{\mathbf{e}_n}{\mathbf{e}_i}$$
    Z definicji bazy ortonormalnej \eqref{eq:def_ortonormal_base}:
    $$\inner{A\mathbf{e}_k}{\mathbf{e}_i} = a_{1k} \times 0 + \ldots + a_{ik}\times 1 + \ldots + a_{nk} \times 0$$
    Więc:
    \begin{equation}\label{eq:def_ortonormal_coeffs}
        a_{ik} = \inner{A\mathbf{e}_k}{\mathbf{e}_i}
    \end{equation}
    Wzór \eqref{eq:def_ortonormal_coeffs} jest wzorem na współczynniki macierzy operatora liniowego w bazie ortonormalnej.
    Niech w tej samej bazie operatorowi $A^*$ odpowiada macierz $(a^*_{ik})$, gdzie $i, k = 1, \ldots, n$.
    Wtedy ze wzoru \eqref{eq:def_ortonormal_coeffs}:
    \begin{equation}\label{eq:local_proof_1}
        a^*_{ik} = \inner{A^*\mathbf{e}_k}{\mathbf{e}_i}
    \end{equation}
    Z definicji operatora sprzężonego \eqref{eq:def_conjugated_operator}:
    $$\inner{A^*\mathbf{e}_k}{\mathbf{e}_i} = \inner{\mathbf{e}_k}{A\mathbf{e}_i}$$
    \begin{equation}\label{eq:local_proof_2}
        \inner{A^*\mathbf{e}_k}{\mathbf{e}_i} = \overline{\inner{A\mathbf{e}_i}{\mathbf{e}_k}}
    \end{equation}
    Z \eqref{eq:def_ortonormal_coeffs}, \eqref{eq:local_proof_1} oraz \eqref{eq:local_proof_2} wynika, że:
    $$a^*_{ik} = \bar{a}_{ki}, \quad (i, k = 1, \ldots, n)$$
    A zatem operatorowi sprzężonemu $A^*$ odpowiada macierz sprzężona.
\end{proof}

\begin{theorem}\label{theorem:properties_of_conjugated_operator}
    Z definicji operatora sprzężonego mamy następujące własności:
    \begin{enumerate}[leftmargin=1cm, itemsep=2pt, label=(\arabic*)]
        \item $(A^*)^* = A$
        \item $(A + B)^* = A^* + B^*$
        \item $(\lambda A)^* = \bar{\lambda}(A^*)$, gdzie $\lambda \in \C$
        \item $(AB)^* = B^*A^*$
    \end{enumerate}
\end{theorem}
\begin{proof}
Dla (1): $(A^*)^* = (a^*_{ik})^* = (\bar{a}_{ki})^* = (a_{ik}) = A$, gdzie $i, k \in \N$

Dla (2):
$(A + B)^* = (a_{ik} + b_{ik})^* = \bar{a}_{ki} + \bar{b}_{ki} = a^*_{ik} + b^*_{ik} = A^* + B^*$, gdzie $i, k \in \N$

Dla (3):
$(\lambda A)^* = (\lambda a_{ik})^* = \overline{(\lambda a_{ki})} = \bar{\lambda}(a^*_{ik}) = \bar{\lambda}(A^*)$, gdzie $\lambda \in \C, i, k \in \N$

Dla (4):
Zakładamy, że możemy wymnożyć macierze $A$ i $B$, a więc liczba kolumn macierzy $A_{(p \times n)}$ zgadza się z liczbą wierszy macierzy $B_{(n \times m)}$. Z tego wynika, że liczba wierszy macierzy $A^*_{(n \times p)}$, zgadza się z liczbą kolumn macierzy $B^*_{(m \times n)}$, gdzie $n, p, m \in \N$.
Element na pozycji $(i, k)$ macierzy $(AB)$ z definicji mnożenia dany jest wzorem: 

$$(AB)_{(i, k)} = \sum_{r=1}^{n}{a_{ir}b_{rk}}$$
Element na pozycji $(i, k)$ macierzy $(AB)^*$ dany jest wzorem:
$$(AB)^*_{(i, k)} = \overline{(AB)}_{(k, i)} = \sum_{r=1}^{n}{\bar{a}_{kr}\bar{b}_{ri}}$$
Dla macierzy $(B^*A^*)$ element $(i, k)$ z definicji mnożenia macierzy jest dany wzorem:
$$(B^*A^*)_{(i, k)} = \sum_{r=1}^{n}{b^*_{ir}a^*_{rk}} = \sum_{r=1}^{n}{\bar{b}_{ri}\bar{a}_{kr}} = \sum_{r=1}^{n}{\bar{a}_{kr}\bar{b}_{ri}}$$
Z tego wynika, że $(B^*A^*)_{(i, k)} = (AB)^*_{(i, k)}$, więc $(AB)^* = B^*A^*$.
\end{proof}

\begin{definition}\label{def:complex_normal_operator}
Operator liniowy $A$ nazywa się \textit{normalnym}, jeśli:
$$AA^* = A^*A$$
\end{definition}
\begin{theorem}
    Operator normalny $A$ w przestrzeni unitarnej $V$ z bazą ortonormalną ma macierz normalną.
\end{theorem}
\begin{proof}
% TODO
\end{proof}

\begin{theorem}\label{theorem:normal_operator_property}
    Niech $A$ operator normalny, wtedy dla dowolnego wektora $v \in V$:
    $$\norm{Av} = \norm{A^*v}$$
\end{theorem}
\begin{proof}
Korzystając z definicji operatora normalnego oraz własności operatora sprzężonego: 
$$\norm{Av}^2 = \inner{Av}{Av} = \inner{A^*Av}{v} = \inner{AA^*v}{v} = \inner{A^*v}{A^*v} = \norm{A^*v}^2$$
Ponieważ $\norm{Av} \ge 0$:
$$\norm{Av} = \norm{A^*v}$$
\end{proof}

\begin{definition}\label{def:h_operator}
Operator liniowy $H$ nazywa się \textit{hermitowskim}, jeśli:
$$H = H^*$$
\end{definition}
\begin{theorem}
    Operator hermitowski $H$ w przestrzeni unitarnej $V$ z bazą ortonormalną ma macierz hermitowską.
\end{theorem}
\begin{proof}
    Z definicji operatora hermitowskiego dla dowolnych dwóch wektorów $\textbf{e}_i$, $\textbf{e}_k$ z bazy ortonormalnej $\base{e}$, gdzie $i, k \in \N$:
    $$\inner{H\textbf{e}_k}{\textbf{e}_i} = \inner{\textbf{e}_k}{H\textbf{e}_i}$$
    Korzystając ze wzoru \eqref{eq:def_ortonormal_coeffs} oraz sprzężonej symetrii iloczynu skalarnego:
    $$ a_{ik} = \inner{H\textbf{e}_k}{\textbf{e}_i} = \inner{\textbf{e}_k}{H\textbf{e}_i} = \overline{\inner{H\textbf{e}_i}{\textbf{e}_k}} = \bar{a}_{ki}, $$ gdzie $i, k \in \{1, \dots, n\}$. Z tego wynika, że operatorowi hermitowskiemu odpowiada macierz hermitowska.
\end{proof}

\begin{definition}\label{def:unitary_operator}
    Operator liniowy $U$ nazywa się \textit{unitarnym}, jeśli dla dowolnych dwóch wektorów $x, y \in V$:
    \begin{equation}\label{eq:def_unitary_operator}
        \inner{Ux}{Uy} = \inner{x}{y}
    \end{equation}
\end{definition}
Równość \eqref{eq:def_unitary_operator} możemy przepisać w następujący sposób na mocy \eqref{eq:def_conjugated_operator}:
$$\inner{Ua}{b} = \inner{a}{U^*b}$$
Podstawiając pod $a = x$, a pod $b = Uy$ otrzymujemy:
$$\inner{Ux}{Uy} = \inner{x}{U^*Uy} = \inner{x}{y}$$
Z tego wynika:
\begin{equation}
    U^*U = I,
\end{equation} gdzie $I$ to operator tożsamościowy (macierz jednostkowa). 

\begin{theorem}
    Operator unitarny $U$ w przestrzeni unitarnej $V$ z bazą ortonormalną ma macierz unitarną.
\end{theorem}
\begin{proof}
    Załóżmy, że $U$ jest operatorem unitarnym w przestrzeni unitarnej z bazą ortonormalną $\base{e}$, i odpowiednią mu macierzą, oznaczymy ją jako $A$.
    Ponieważ $\base{e}$ jest bazą ortonormalną, to $U(\mathbf{e}_1), \ldots, U(\mathbf{e}_n)$ też jest bazą ortonormalną, dlatego że $U$ zachowuje iloczyn skalarny tych wektorów. Te wektory traktowane jako wektory kolumnowe tworzą macierz unitarną, tworząc jednocześnie macierz $A$.

    Na odwrót, załóżmy, że macierz $A$ jest unitarna. Niech $x, y \in V$. Wówczas, korzystając z tw. \ref{theorem:properties_of_conjugated_operator} (4 własność):
    $$\inner{U(x)}{U(y)} = \inner{Ax}{Ay} = (Ay)^*(Ax) = x^*A^*Ay = x^*y = \inner{x}{y},$$
    a to oznacza, że $U$ jest operatorem unitarnym.
\end{proof}

Operatory hermitowskie oraz unitarne są szczególnymi przypadkami operatora normalnego.
W przestrzeni euklidesowej tym operatorom odpowiadają operatory symetryczne oraz ortogonalne. Rozpatrzymy to w następnym podrozdziale.
